\section{خوارزمية التعرف على الكلام باستخدام مكتبة Vosk}
تم اعتماد حزمة \textenglish{Vosk} كأداة أساسية للتعرف على الكلام في هذا المشروع بناءً على ميزاتها التقنية المتقدمة التي تلائم الأنظمة المدمجة والمستقلة:

\begin{itemize}
    \item \textbf{العمل دون اتصال بالإنترنت (\textenglish{Offline Processing}):} لا تتطلب الحزمة أي اتصال بخوادم سحابية خارجية، مما يضمن أمان البيانات واستمرارية تشغيل الذراع الروبوتية في بيئات العمل المعزولة أو الصناعية التي تفتقر إلى شبكة إنترنت مستقرة.
    \item \textbf{الاستهلاك المنخفض للموارد:} توفر الحزمة نماذج لغوية مصغرة (\textenglish{Small Language Models}) مصممة خصيصاً للعمل بكفاءة واستهلاك معتدل للذاكرة على المعالجات المدمجة ذات الموارد المحدودة.
    \item \textbf{الدعم المتعدد للغات المتنوعة:} تتيح الحزمة إمكانية دمج نماذج صوتية مدربة تدريباً عالياً للتعرف على الكلمات بدقة ممتازة ونسبة خطأ ضئيلة حتى في وجود ضجيج محيطي مصاحب لبيئة التشغيل.
\end{itemize}

تعتمد آلية عمل خوارزمية \textenglish{Vosk} في تفسير الأوامر على دمج نموذجين رياضيين أساسيين:
\begin{enumerate}
    \item \textbf{النموذج الصوتي (\textenglish{Acoustic Model}):} يقوم باستقبال الإشارة الصوتية الرقمية الملتقطة، ثم يعمل على استخلاص ميزاتها الرياضية والترددية عبر حساب معاملات الـ \textenglish{Mel-Frequency Cepstral Coefficients (MFCC)}، وتحويلها إلى احتمالات إحصائية تمثل الأصوات اللغوية الأساسية البسيطة (\textenglish{Phonemes}).
    \item \textbf{النموذج اللغوي (\textenglish{Language Model}):} يمثل البنية القواعدية والمعجمية للكلمات، حيث يقوم بربط الأصوات المستخلصة من النموذج الصوتي وتجميعها لتشكيل الكلمات والنصوص النهائية بناءً على مصفوفة احتمالية وشبكة بحث مسبقة الصنع للتنبؤ بالكلمة التالية.
\end{enumerate}